\documentclass[11pt]{article}

\usepackage[T1]{fontenc}
\usepackage[utf8]{inputenc}
\usepackage[spanish,es-tabla]{babel}
\usepackage{amsmath,amssymb}
\usepackage{booktabs}
\usepackage{geometry}
\geometry{margin=2.4cm}
\usepackage[hidelinks]{hyperref}
\usepackage{xcolor}
\newcommand{\estado}[1]{\textcolor{blue!60!black}{\textbf{[#1]}}}
\newcommand{\ok}{\textcolor{green!50!black}{\checkmark}}
\newcommand{\falla}{\textcolor{red!70!black}{$\times$}}
\newcommand{\kk}{\Bbbk}
\newcommand{\kQ}{\kk Q}

\title{Bit\'acora de experimentos --- camera-ready CIARP\\
\large \emph{Path Algebras of Quivers as a Sparse Multi-Hop Attention
for Mitigating Over-Squashing in GNNs} (env\'io 15, \textbf{aceptado})}
\author{C. I. Zainea Maya \and A. Moreno Ca\~nadas}
\date{Iniciada: 26 de agosto de 2026 --- \'ultima actualizaci\'on: 26-08-2026 (cierre de fase)}

\begin{document}
\maketitle

\section{Contexto y observaciones de los jurados}

El art\'iculo fue \textbf{aceptado} en CIARP (Revisor 1: \emph{accept},
puntajes 4--5 en todos los criterios; Revisor 2: \emph{weak reject}, 3--4).
Esta bit\'acora registra la fase experimental posterior a la aceptaci\'on y el
impacto de cada resultado sobre las observaciones. Regla de trabajo: el
manuscrito \textbf{no se modifica} hasta cerrar esta fase.

\begin{description}
  \item[R1.1] Discutir el costo de precomputar $A^k$ en grafos masivos.
  \item[R1.2] Hipotetizar un escenario concreto donde el ideal $I$ agregue
  valor (p.\,ej.\ clases de caminos $\to$ cabezas de atenci\'on).
  \item[R1.3] Validar en un dataset real con over-squashing severo.
  \item[R1.4] Restaurar nombres en las referencias.
  \item[R2] ``\emph{Much of the algebraic machinery is not essential to the
  method}'': el modelo usar\'ia solo el soporte de $A^k$ como m\'ascara
  dispersa; las ganancias fuertes son sint\'eticas; en Peptides-func empata
  con GAT y queda bajo los baselines afinados de la literatura.
\end{description}

\paragraph{Estrategia experimental.} R2 tiene raz\'on en que \emph{el modelo}
solo consume el soporte. La respuesta no es discutirlo sino mostrar que el
\'algebra hace un segundo trabajo que la m\'ascara no puede hacer:
\textbf{medir} el over-squashing y \textbf{predecir} el r\'egimen donde Walk
Attention gana. Tres frentes: (1) un \'indice algebraico de severidad
validado contra los resultados ya publicados; (A) densificar esa curva con
nuevas corridas; (B) un benchmark real elegido \emph{por} el \'indice.

\section{Acciones de texto ya aplicadas (commit \texttt{b29769b}, 26-08)}

\begin{center}\small
\begin{tabular}{p{3.6cm}p{9.4cm}l}
\toprule
\textbf{Observaci\'on} & \textbf{Acci\'on} & \textbf{Estado}\\
\midrule
R1.4 & Autores restaurados; entradas bib reales
(\texttt{zainea2026aiq}, repos de c\'odigo); DOI de CAsimulations
\texttt{10.15332/23393076.11217}. & \estado{hecho}\\
R1.1 & \S4.1: precomputaci\'on por $k{-}1$ productos dispersos, costo por
no-ceros intermedios; truncado/muestreo en grafos masivos. & \estado{hecho}\\
R1.2 & Trabajo futuro: grafo heterog\'eneo, ideal que agrupa composiciones
equivalentes de relaciones tipadas $\to$ cabezas compartidas. & \estado{hecho}\\
Menores & Fig.~2 al margen; Prop.~5 con $\beta/\sqrt{C}$ y valores base
distintos; \textsc{RingTransfer} acreditado a Bodnar et al.\ 2021; tablas ES
traducidas. & \estado{hecho}\\
R1.3, R2 & Frentes experimentales de esta bit\'acora. & \estado{en curso}\\
\bottomrule
\end{tabular}
\end{center}

% =====================================================================
\section{Experimento 1 --- el \'indice de severidad $\Psi$ y la curva
dosis-respuesta}
% =====================================================================

\subsection{Derivaci\'on del \'indice}

El argumento de capacidad del paper (Prop.~3--4) es: tras $k$ capas de paso
de mensajes, \emph{todas} las cadenas de longitud $k$ que entran al nodo $j$
--- una por cada elemento de la base de $\bigoplus_i e_j\kQ_k e_i$ --- deben
codificarse en un vector de ancho $w$. Si las cadenas fueran distinguibles y
equiprobables, distinguirlas requiere $\log_2\bigl(\sum_i(A^k)_{ij}\bigr)$
bits, mientras el embedding ofrece del orden de $\log_2 w$. El \'indice es el
exceso:
\[
\Psi_k(j)\;=\;\log_2\Bigl(\textstyle\sum_i (A^k)_{ij}\Bigr)\;-\;\log_2 w,
\qquad
\sum_i (A^k)_{ij}=\sum_i \dim_\kk\bigl(e_j\,\kQ_k\,e_i\bigr).
\]
Es \'algebra graduada pura: dimensiones de las componentes $e_j\kQ_k e_i$
(Prop.~1 del paper), agregadas por objetivo. Es la versi\'on por objetivo del
diagn\'ostico por pares $H_k(i,j)=\log_2(A^k)_{ij}$ ya implementado en
\texttt{aiq.gnn.over\_squashing\_diagnostic}. Predicci\'on cualitativa: con
$\Psi\le0$ el paso de mensajes tiene capacidad suficiente y la atenci\'on
multi-salto no deber\'ia aportar; con $\Psi\gg0$ el colapso es inevitable y
Walk Attention (que lee las fuentes \emph{directamente}, sin atravesar el
canal comprimido) deber\'ia dominar.

\subsection{Protocolo}

Script \texttt{run\_severity\_index.py}. Para cada benchmark ya publicado se
computa $\Psi$ en \emph{el objetivo y alcance de decisi\'on} --- el nodo que
porta la etiqueta y el alcance $k$ que la resuelve --- y se cruza con la
ventaja medida $=$ media WA $-$ mejor baseline de un salto, de los CSV por
semilla de \texttt{results/tables/}:
\begin{itemize}
  \item \textbf{Bottleneck} ($K{=}5$, $M{=}4$, $w{=}16$): objetivo final,
  $k=d{+}1$; masa entrante $=K\!\cdot\!M^{d}$, luego
  $\Psi=\log_2(5\cdot4^{d})-4$.
  \item \textbf{RingTransfer} ($w{=}32$): objetivo antipodal, $k=n/2$; el
  ciclo no dirigido tiene masa entrante $2^{n/2}$, luego $\Psi=n/2-5$.
  \item \textbf{Peptides-func} ($w{=}80$, alcances $1$--$4$): por mol\'ecula,
  m\'aximo sobre objetivos y alcances; promedio sobre las 10\,873 mol\'eculas
  del split de entrenamiento.
\end{itemize}
Sin ninguna corrida nueva: el experimento solo reinterpreta resultados ya
publicados a la luz del \'indice.

\subsection{Resultados}

\begin{center}
\begin{tabular}{lrrrr}
\toprule
\textbf{Benchmark} & $\Psi$ (bits) & \textbf{alcance} & \textbf{WA} &
\textbf{ventaja}\\
\midrule
ring $n{=}6$        & $-2.00$ & 3   & 1.000 & 0.000\\
Peptides-func       & $-0.75$ & 1--4& 0.526 & 0.000\\
ring $n{=}10$       & $0.00$  & 5   & 1.000 & 0.000\\
ring $n{=}14$       & $+2.00$ & 7   & 1.000 & $+0.189$\\
bottleneck $d{=}2$  & $+2.32$ & 3   & 1.000 & $+0.550$\\
ring $n{=}18$       & $+4.00$ & 9   & 1.000 & $+0.466$\\
bottleneck $d{=}3$  & $+4.32$ & 4   & 1.000 & $+0.708$\\
\bottomrule
\end{tabular}
\end{center}

Spearman $\rho=0.927$ ($p=0.0027$); Pearson $r=0.901$ ($p=0.0056$); $n=7$.
Salida: \texttt{results/tables/severity\_index.csv}.

\subsection{Lectura}

\begin{enumerate}
  \item \textbf{Umbral, no solo tendencia.} Los tres casos con $\Psi\le0$
  tienen ventaja \emph{exactamente} $0.000$; la ventaja aparece al cruzar
  $\Psi=0$ y crece con $\Psi$. El corte est\'a donde el argumento de
  capacidad lo pone, sin par\'ametros ajustados.
  \item \textbf{El empate de Peptides queda predicho.} $\Psi=-0.75<0$: en
  esas mol\'eculas ($w{=}80$, grado t\'ipico $2$--$3$, alcances $\le4$) la
  masa de recorridos entrante a\'un cabe en el embedding. El caso que R2 usa
  como evidencia en contra es un punto \emph{sobre la curva} que el \'indice
  anticipa sin entrenar.
  \item \textbf{Impacto sobre R2.} La ``maquinaria prescindible'' (las
  dimensiones graduadas con multiplicidad --- no el patr\'on binario de
  $A^k$, que es el mismo para masa $2$ o $22\,000$) es la que produce
  $\Psi$. La m\'ascara no mide; el \'algebra s\'i.
  \item \textbf{Limitaci\'on.} $n=7$ puntos de 3 familias distintas mezcla
  tareas y anchos; la comparaci\'on justa es \emph{dentro} de cada familia
  (mon\'otona en ambas) y el Frente A densifica.
\end{enumerate}

% =====================================================================
\section{Frente A --- densificaci\'on de la curva (resultados)}
% =====================================================================

\estado{completo} Lanzado 26-08 08:24, cerrado 08:55; mismo protocolo
estabilizado del paper (5 semillas, AdamW, warmup, coseno, clip; LN en WA).
Scripts \texttt{run\_ring\_fillin.py} ($n\in\{12,16,20\}$),
\texttt{run\_bottleneck\_fillin.py} ($d{=}4$).

\subsection{Resultados por configuraci\'on (media $\pm$ desv., 5 semillas)}

\begin{center}
\begin{tabular}{lrrrrr}
\toprule
\textbf{Config.} & $\Psi$ & \textbf{GCN} & \textbf{GAT} & \textbf{WA} &
\textbf{ventaja}\\
\midrule
ring $n{=}12$ & $1.00$ & $0.966\pm0.069$ & $1.000\pm0.000$ &
$1.000\pm0.000$ & $0.000$\\
ring $n{=}16$ & $3.00$ & $0.532\pm0.274$ & $0.459\pm0.229$ &
$1.000\pm0.000$ & $+0.468$\\
ring $n{=}20$ & $5.00$ & $0.418\pm0.289$ & $0.221\pm0.008$ &
$0.894\pm0.156$ & $+0.476$\\
bottleneck $d{=}4$ & $6.32$ & (azar, Prop.~4) & $0.366\pm0.174$ &
$1.000\pm0.000$ & $+0.634$\\
\bottomrule
\end{tabular}
\end{center}

(En $d{=}4$ el operador de recorrido fijo dio $0.571\pm0.050$, entre GAT y
WA, como en $d{=}2,3$: cruza el cuello pero no selecciona.)

\subsection{Contraste con las predicciones pre-registradas}

\begin{center}\small
\begin{tabular}{p{4.5cm}p{5.4cm}l}
\toprule
\textbf{Predicci\'on (antes de correr)} & \textbf{Resultado} & \\
\midrule
$n{=}12$ ($\Psi{=}1$): ventaja peque\~na pero positiva &
ventaja $=0.000$ (GAT a\'un resuelve exacto) & \falla\\
$n{=}16$ ($\Psi{=}3$): ventaja entre las de $n{=}14$ y $n{=}18$ &
$+0.468$, entre $+0.189$ y $+0.466$ (borde superior) & \ok\\
$n{=}20$ ($\Psi{=}5$): ventaja $\ge$ la de $n{=}18$ & $+0.476$ vs
$+0.466$ \ok; pero WA ya no es exacta a 10 capas ($0.894$; semillas
$\{1,1,.82,.65,1\}$) & \ok$\sim$\\
$d{=}4$ ($\Psi{=}6.32$): ventaja $\ge+0.7$, WA $=1.000$ &
WA $=1.000$ \ok; ventaja $+0.634$ (GAT vari\'o alto) & $\sim$\\
\bottomrule
\end{tabular}
\end{center}

La predicci\'on fallida es informativa: \textbf{el umbral efectivo no est\'a
en $\Psi=0$ sino entre $\Psi=1$ y $\Psi=2$}. Lectura mec\'anica: un bit de
exceso a\'un puede absorberse con codificaci\'on no lineal del embedding
(la cota $\log_2 w$ es conservadora); con dos o m\'as bits el colapso es
sistem\'atico. El \'indice queda entonces con una zona de operaci\'on:
$\Psi\le0$ seguro-sin-ventaja, $\Psi\ge2$ seguro-con-ventaja, $0<\Psi<2$
zona de transici\'on.

\subsection{El fen\'omeno de transici\'on: bimodalidad por semilla}

Los promedios ocultan la estructura m\'as reveladora. Por semilla, en
$n{=}16$: GCN $\{0.61, 0.21, 0.21, 0.81, 0.82\}$, GAT
$\{0.62, 0.21, 0.22, 0.80, 0.46\}$. Los baselines de un salto \textbf{no
degradan gradualmente: o resuelven ($\approx0.8$--$1.0$) o caen al azar
($\approx0.21$)}, y la proporci\'on de semillas que caen crece con $\Psi$
hasta que en $n{=}20$/$d{=}4$ casi todas caen. La desviaci\'on est\'andar
por eso \emph{sube} en la transici\'on ($\pm0.27$ en $n{=}16$;
$\pm0.17$ en $d{=}4$) y colapsar\'ia a $0$ en ambos extremos. Walk
Attention, en cambio, da $1.000\pm0.000$ en \textbf{todas} las
configuraciones y semillas de todos los benchmarks sint\'eticos --- la
estabilidad entre semillas es en s\'i misma un resultado: el soporte
algebraico elimina la loter\'ia de inicializaci\'on que decide si el canal
comprimido logra o no codificar la se\~nal.

\subsection{Curva final}

Con los 11 puntos: $\Psi\in\{-2,-0.75,0,1\}\mapsto$ ventaja $0.000$;
$\Psi\in\{2,2.32,3,4,4.32,5,6.32\}\mapsto
\{+0.19,+0.55,+0.47,+0.47,+0.71,+0.48,+0.63\}$.
\textbf{Spearman $\rho=0.875$ ($p=4.3\times10^{-4}$), Pearson $r=0.880$.}
Plana en cero hasta $\Psi=1$, salto en $\Psi=2$, meseta alta con ruido
dominado por la bimodalidad del mejor baseline. Nota honesta: en $n{=}20$
(10 capas) la propia WA pierde exactitud en 2 de 5 semillas --- la
profundidad de la pila de atenci\'on introduce su propia dificultad de
optimizaci\'on, independiente del squashing.

% =====================================================================
\section{Experimento 2 --- cribado de redes reales con $\Psi$}
% =====================================================================

\subsection{Protocolo}

Script \texttt{run\_screen\_real.py}. La masa entrante se computa como
$\mathbf{1}^{T}A^{k}$ por iteraci\'on vector--matriz, $O(k\cdot\mathrm{nnz})$
--- sin materializar $A^k$: el mecanismo de escalabilidad respondido a R1.1,
funcionando sobre 421K aristas en segundos. Direcciones: \emph{as-is}
(citante$\to$citado: las cadenas de citaci\'on se acumulan en los papers
citados) y \emph{reversa}. Donde las potencias completas caben se cuenta
adem\'as el sustrato del cociente: pares con $(A^k)_{ij}>1$, i.e.\
$\dim(e_j\kQ_k e_i)>1$ (recorridos paralelos).

\subsection{Resultados (selecci\'on; tabla completa en
\texttt{real\_network\_severity\_screen.csv})}

\begin{center}\small
\begin{tabular}{llrrrrr}
\toprule
\textbf{Red} & \textbf{dir.} & $k$ & \textbf{m\'ax bits} &
\textbf{frac.}$>w{=}64$ & \textbf{pares paralelos} & \textbf{m\'ax mult.}\\
\midrule
cit-HepPh & reversa & 2 & 13.2 & 0.51 & $1.1\times10^{6}$ & 110\\
(34\,546 n.) & reversa & 3 & 17.4 & 0.68 & $9.4\times10^{6}$ & 1\,481\\
 & reversa & 4 & 21.4 & 0.72 & $3.4\times10^{7}$ & \textbf{22\,702}\\
Cora (2\,708 n.) & reversa & 6 & 9.5 & 0.22 & 14\,711 & 125\\
 & reversa & 8 & 11.8 & 0.44 & 21\,933 & 772\\
Ca\~nadas (115 n.) & reversa & 3 & 6.1 & 0.01 & 37 & 6\\
\bottomrule
\end{tabular}
\end{center}

\subsection{Lectura}

\begin{itemize}
  \item \textbf{cit-HepPh es el r\'egimen severo real (R1.3), con margen
  enorme}: $\Psi\approx+15$ bits sobre $w{=}64$ a $k{=}4$; el sint\'etico
  m\'as severo del paper tiene $+4.3$. No es un caso constru\'ido: es la
  estructura natural de una red de citaci\'on con hubs.
  \item \textbf{La multiplicidad extrema replantea el cociente (R1.2).} El
  resultado negativo del paper (colapsar 2 recorridos paralelos destruye
  se\~nal) se obtuvo con multiplicidades de decenas. Aqu\'i hay pares con
  22\,702 cadenas paralelas: la hip\'otesis natural se invierte ---
  \emph{organizar} esas clases (cociente parcial $\to$ cabezas) deber\'ia
  ganar sobre ignorarlas y sobre colapsarlas. Es el experimento del espectro
  de ideales del Frente B.
  \item \textbf{Cora} es punto intermedio; \textbf{Ca\~nadas} (leve, mult.\
  $\le6$) sirve como ejemplo trabajado, no como benchmark.
\end{itemize}

% =====================================================================
\section{Frente B --- benchmark real cit-HepPh}
% =====================================================================

\subsection{El problema real}

\textbf{Detecci\'on temprana de art\'iculos influyentes}: en 1997, con el
grafo de citaciones disponible, se\~nalar cu\'ales papers de 1996--97
ser\'an los m\'as citados en 1998--2000. El conteo de citas es un indicador
retrasado; el predictor barato (apego preferencial: citas presentes
predicen futuras) falla justo en los casos valiosos --- los que despegan
tarde. La se\~nal correctora es estructural: \emph{qui\'en} te cita (¿papers
que a su vez ser\'an influyentes?) y por cu\'antas cadenas independientes
llega tu influencia --- informaci\'on que vive en $e_j\kQ_k e_i$ con su
multiplicidad, a profundidades $2$--$4$.

\subsection{Construcci\'on del dataset (\texttt{run\_hepph\_dataset.py})}

Corte 31-12-1997. Grafo de contexto: 13\,745 papers fechados y 98\,307
citaciones entre ellos. Cohorte: 6\,635 papers de 1996--97. Futuro: 51\,632
citaciones recibidas en 1998--2000 (mediana 3, m\'aximo 298). Features
m\'inimos e id\'enticos para todos los modelos: a\~no de publicaci\'on,
$\log(1{+}\mathrm{in})$, $\log(1{+}\mathrm{out})$ al corte, estandarizados.
Ids cruzados de SNAP (\texttt{11<id>}) normalizados. Severidad verificada
\emph{en el subgrafo de contexto}: m\'ax $18.4$ bits a $k{=}4$
($\Psi=+12.4$ sobre $w{=}64$), multiplicidad m\'axima 4\,724,
$\mathrm{nnz}(A^4)=2.7$M.

\subsection{Cinco variables objetivo}

\begin{center}\small
\begin{tabular}{llrr}
\toprule
 & \textbf{Definici\'on (binaria)} & \textbf{balance} & \textbf{definidos}\\
\midrule
T1 impacto estrat. & top-cuartil de citas futuras \emph{dentro} de su
estrato de in-degree $\{0,\,1$--$2,\,3$--$9,\,10^{+}\}$ & 0.23 & 6\,635\\
T2 early riser & top-cuartil del residuo de $\log(1{+}c_{\text{fut}})$
sobre grado y edad & 0.25 & 6\,635\\
T3 amplitud & fracci\'on de pares de citantes futuros \emph{independientes}
(sin citaci\'on mutua ni referencia compartida) sobre mediana & 0.50 &
4\,187\\
T4 persistencia & fracci\'on de citas futuras tard\'ias ($\ge$ jul-1999)
sobre mediana & 0.48 & 3\,578\\
T5 alcance & top-cuartil del crecimiento de la 2-vecindad de impacto
$|A_2(c)|$ & 0.25 & 6\,635\\
\bottomrule
\end{tabular}
\end{center}

T1/T2 neutralizan el apego preferencial por construcci\'on (estratificado /
residuo). T3 es el objetivo cuya \emph{definici\'on} es multiplicidad de
cadenas independientes --- el terreno del cociente. T5 usa el grado de
impacto de la tesis como variable objetivo: predecir una cantidad
algebraica futura desde el \'algebra presente.

\subsection{Modelos y protocolo (\texttt{run\_hepph\_models.py})}

Transductivo sobre el grafo dirigido (citante$\to$citado, los mensajes
fluyen de citantes a citados); divisi\'on estratificada 60/20/20 de la
cohorte; 5 semillas; p\'erdida ponderada por clase; m\'etrica: exactitud
balanceada en test. Modelos: MLP (piso sin grafo), GCN, GAT (4 cabezas),
Walk Attention (una capa por alcance $k=1..4$, soporte truncado top-32
fuentes por multiplicidad --- ``Scaling the support'' operativo:
$\mathrm{nnz}\le n\cdot32$ por alcance). Todos con ancho 64 y la receta
estabilizada del paper. Smoke test (3 \'epocas, 2 capas):
mlp $0.50$ / gcn $0.54$ / gat $0.56$ / wa $0.61$ --- ordenamiento
alentador, no concluyente. \estado{corrida completa lanzada 26-08;
log en \texttt{results/logs/hepph\_models.log}}

\subsection{Hip\'otesis pre-registradas (antes de ver resultados)}

\begin{enumerate}
  \item En T1/T2, orden esperado MLP $<$ GCN/GAT $<$ WA; la ganancia de WA
  sobre GAT ser\'a mayor en los estratos de grado bajo (donde el conteo no
  informa y solo la estructura profunda distingue).
  \item La brecha m\'axima WA--GAT aparecer\'a en \textbf{T3} (amplitud), el
  objetivo definido por multiplicidad de cadenas.
  \item En T5 todos los modelos con grafo superar\'an ampliamente al MLP
  (el objetivo es estructural), y WA $>$ GAT por el alcance.
  \item (Fase 2, espectro de cocientes sobre T3:) $I_{\mathrm{hub}}\gtrsim
  I_0\gg I_{\mathrm{full}}$ con multiplicidad extrema --- el primer
  resultado \emph{positivo} del cociente.
\end{enumerate}

\subsection{Resultados (100 corridas base $+$ 50 de control)}

\begin{center}
\begin{tabular}{lrrrrr}
\toprule
\textbf{Modelo} & \textbf{T1} & \textbf{T2} & \textbf{T3} & \textbf{T4} &
\textbf{T5}\\
\midrule
MLP (sin grafo)      & 0.672 & 0.565 & 0.527 & 0.519 & 0.818\\
GCN (sin residual)   & 0.636 & 0.548 & 0.542 & 0.522 & 0.744\\
GAT (sin residual)   & 0.666 & 0.546 & 0.546 & 0.523 & 0.785\\
GCN $+$ residual     & 0.664 & 0.570 & 0.545 & 0.522 & 0.811\\
GAT $+$ residual     & 0.673 & 0.584 & 0.550 & 0.529 & 0.810\\
\textbf{Walk Attention} & \textbf{0.673} & \textbf{0.594} & \textbf{0.559}
& 0.526 & \textbf{0.828}\\
\bottomrule
\end{tabular}
\end{center}
(Exactitud balanceada en test, media de 5 semillas; desviaciones
$0.004$--$0.028$. Archivo \texttt{hepph\_models\_5seeds.csv}.)

\subsection{El confound y su control}

La primera lectura (GCN/GAT \emph{por debajo} del MLP en T1, T2, T5:
``la agregaci\'on local da\~na'') era demasiado buena. WA lleva t\'ermino
residual $W_{\text{root}}$ de f\'abrica; los GCN/GAT del runner no llevaban
skip connections. El control (\texttt{gcn\_res}, \texttt{gat\_res}, mismas
semillas) lo confirma: \textbf{con residual, GCN/GAT recuperan el nivel del
MLP y m\'as}; el efecto ``la agregaci\'on local da\~na'' era en su mayor
parte el confound. Se descarta esa lectura.

\subsection{Lo que sobrevive al control}

WA sigue siendo el mejor modelo en T1 (empate), T2, T3 y T5, con
ventajas peque\~nas por objetivo ($+0.009$ a $+0.018$ sobre GAT$+$res)
pero \textbf{consistentes}: prueba pareada por (objetivo, semilla),
excluyendo el nulo T4 ($n=20$ pares):
\begin{center}
\begin{tabular}{lrrr}
\toprule
 & \textbf{media} & \textbf{pares positivos} & \textbf{Wilcoxon $p$}\\
\midrule
WA $-$ GAT$+$res & $+0.009$ & 15/20 & 0.004\\
WA $-$ GCN$+$res & $+0.016$ & 18/20 & 0.0003\\
WA $-$ MLP       & $+0.018$ & 16/20 & 0.0009\\
\bottomrule
\end{tabular}
\end{center}

\subsection{Contraste con las hip\'otesis y lectura revisada}

\begin{itemize}
  \item H1 (orden MLP $<$ GCN/GAT $<$ WA en T1/T2): \falla\ en T1 (empate
  general: los features de grado agotan la se\~nal); \ok\ en T2 con
  m\'argenes peque\~nos.
  \item H2 (brecha m\'axima en T3): \falla\ --- la brecha m\'axima est\'a
  en T2 (residuo ortogonal al grado), T3 segunda.
  \item H3 (grafo $\gg$ MLP en T5): \falla\ --- el grado actual predice el
  crecimiento de la 2-vecindad; solo WA supera al MLP, por poco.
  \item T4: nulo para todos.
\end{itemize}

\textbf{Refinamiento del \'indice.} $\Psi\approx+12$ predec\'ia una ventaja
\emph{decisiva} y se observ\'o una ventaja \emph{peque\~na y consistente}.
La explicaci\'on es una distinci\'on que el Exp.~1 no pod\'ia revelar
porque sus etiquetas fueron dise\~nadas para depender de la informaci\'on
profunda: $\Psi$ mide cu\'anta informaci\'on de recorridos \emph{se pierde}
en el canal, no cu\'anta de esa informaci\'on \emph{necesita la tarea}. Es
condici\'on necesaria, no suficiente: con $\Psi\le1$ no hay ventaja
posible (Peptides, anillos cortos); con $\Psi\gg0$ la ventaja es posible y
su magnitud la fija la dependencia del label respecto de la estructura
profunda --- alta en las tareas sint\'eticas de recuperaci\'on, modesta en
la predicci\'on de impacto con features de grado (donde el grado ya resume
casi todo lo predecible). Esta es la formulaci\'on correcta para el paper.

\subsection{Decisi\'on sobre las fases 2--3}

El espectro de cocientes sobre HepPh (fase 2) medir\'ia diferencias del
orden de $0.01$, en el l\'imite de resoluci\'on de 5 semillas; se pospone,
salvo que se aumente a 10 semillas. $I_{\mathrm{ring}}$ en Peptides (fase 3)
sigue siendo la pregunta abierta del paper y se mantiene en cola.

% =====================================================================
\section{Cierre de fase: qu\'e queda establecido}
% =====================================================================

\begin{enumerate}
  \item \textbf{Establecido, fuerte.} El \'indice $\Psi$ y la curva
  dosis-respuesta (11 puntos, $\rho=0.875$, umbral en $(1,2]$ bits, 4 ceros
  exactos incluido Peptides). Responde directamente a R2: el \'algebra
  graduada con multiplicidad predice el r\'egimen; la m\'ascara binaria no
  puede.
  \item \textbf{Establecido, moderado.} En una red real de severidad
  extrema (cit-HepPh, $\Psi\approx+12$), WA es el mejor modelo en 4 de 5
  objetivos con ventaja peque\~na pero significativa ($p=0.004$ pareado)
  sobre GAT con residual, con soporte truncado top-32 (escalabilidad
  operativa). Responde a R1.3 con honestidad: ventaja real, no dram\'atica.
  \item \textbf{Aprendido.} $\Psi$ es necesario, no suficiente; la
  magnitud de la ventaja depende de cu\'anto el label necesita la
  informaci\'on profunda. Y: los controles de arquitectura (residuales)
  son obligatorios antes de leer un ``la agregaci\'on local da\~na''.
  \item \textbf{Pendiente.} $I_{\mathrm{ring}}$ en Peptides; espectro de
  cocientes con m\'as semillas.
\end{enumerate}

\paragraph{Propuesta para el camera-ready (a decidir).} Una subsecci\'on
breve tras los experimentos: definici\'on de $\Psi$ (3 l\'ineas), la curva
como tabla compacta o figura, y un p\'arrafo con cit-HepPh (T2/T3/T5 y el
test pareado). Esto convierte ``regime-dependent'' en una predicci\'on
cuantitativa y a\~nade el dataset real; cabe en $\sim$2/3 de p\'agina si se
comprime el ejemplo de la Secci\'on~3.

% =====================================================================
\section{Reproducibilidad}
% =====================================================================

\begin{itemize}
  \item Repositorio en \texttt{b29769b} (rama \texttt{main}) $+$ scripts
  nuevos sin comprometer: \texttt{run\_severity\_index.py},
  \texttt{run\_screen\_real.py}, \texttt{run\_ring\_fillin.py},
  \texttt{run\_bottleneck\_fillin.py}, \texttt{run\_hepph\_dataset.py},
  \texttt{run\_hepph\_models.py}.
  \item Salidas: \texttt{results/tables/severity\_index.csv},
  \texttt{real\_network\_severity\_screen.csv},
  \texttt{bottleneck\_fillin\_5seeds.csv},
  \texttt{ring\_transfer\_fillin\_5seeds.csv},
  \texttt{hepph\_models\_5seeds.csv} (150 filas: base $+$ control);
  dataset en \texttt{data/hepph/}; logs en \texttt{results/logs/}.
  \item Entorno: venv \texttt{oversquash} --- Python 3.12.10, PyTorch
  2.13.0+cpu, PyG 2.8.0.post1, NumPy 2.5.2, SciPy 1.18.1;
  \texttt{oversquash} y \texttt{aiq-quivers} editables.
  \item Datos crudos: sint\'eticos por semilla; Peptides-func en
  \texttt{data/lrgb/}; cit-HepPh/Cora/Ca\~nadas en
  \texttt{../quivers\_analysis/data/}.
\end{itemize}

\end{document}
